%%%
%
% $Autor: Wings $
% $Datum: 2021-05-14 $
% $Pfad: GitLab/CornerBLending $
% $Dateiname: Hints
% $Version: 4620 $
%
% !TeX spellcheck = de_DE/GB
%
%%%



\chapter{Hints}

\begin{tcolorbox}[colback=blue!5!white,
	colframe=blue!60!black,
	title={\faHandshake\hspace{0.5em}Welcome to Your Smart License Detection System},
	fonttitle=\bfseries\large,
	arc=3mm,
	boxrule=1pt]
	
	\noindent
	\textbf{Dear User,}\\[4pt]
	Welcome to the \textbf{Automatic License Plate Detection and Recognition System}, designed and developed on the powerful \textbf{NVIDIA Jetson Nano (reComputer J1010)} platform.  
	This system represents a fusion of cutting-edge edge AI and embedded vision—offering real-time intelligence for smart parking, traffic automation, and access management.
	
	\vspace{0.5em}
	The device automatically captures live video through a \textbf{Microsoft LifeCam Show} camera, detects vehicle license plates using the \textbf{YOLOv8 deep learning model}, and accurately reads plate characters through \textbf{EasyOCR}.  
	All results are securely stored in a local \textbf{SQLite database}, including timestamps and detection confidence, enabling offline operation and audit-friendly logging.
	
	\vspace{0.5em}
	The system is built for real-world reliability:
	\begin{itemize}
		\item Works fully offline—no cloud connection required.
		\item Optimized for the Jetson Nano’s 128-core GPU and compact aluminum case.
		\item Integrates seamlessly into automated gates, parking systems, or campus entry points.
	\end{itemize}
	
	\vspace{0.5em}
	\noindent
	Thank you for choosing this innovation in edge AI automation.  
	We wish you a smooth and inspiring experience with your \textbf{License Plate Detection System}.
	
	\vspace{0.3cm}
	\noindent
	\textit{— The Project Team, Hochschule Emden-Leer}
\end{tcolorbox}

% --------------------------------------------------
% Section 1: Safety Guidelines
% --------------------------------------------------
\section{Pre Safety Guidelines}

Safety is an essential aspect of setting up and operating the \textbf{License Plate Detection System} on the NVIDIA Jetson Nano platform.  
The following guidelines ensure safe operation, hardware longevity, and data protection.

\subsection{Electrical Safety}
\begin{itemize}
	\item Always use the original or a certified \textbf{5V/4A USB-C power adapter} to avoid undervoltage or overcurrent damage.
	\item Ensure that all power connections are properly inserted before switching on the unit.
	\item Do not power the Jetson Nano directly from unstable USB ports or laptop chargers.
	\item Disconnect power before connecting or disconnecting the USB camera, HDMI cable, or other peripherals.
	\item Avoid operating the system in wet or humid environments to prevent electrical short circuits.
\end{itemize}

\subsection{Thermal and Environmental Safety}
\begin{itemize}
	\item The reComputer J1010 enclosure provides \textbf{passive cooling} through its aluminum case; do not cover ventilation areas during operation.
	\item Maintain ambient room temperature below \textbf{40°C} for optimal performance.
	\item Avoid placing the system near direct sunlight, heaters, or enclosed cabinets without airflow.
	\item If prolonged high workloads are expected (e.g., continuous YOLO inference), consider adding a small external fan.
\end{itemize}

\subsection{Operational Safety}
\begin{itemize}
	\item Place the device on a \textbf{flat, stable, and non-conductive} surface to prevent vibration or accidental falls.
	\item Handle the Jetson Nano carefully to avoid static discharge; use anti-static wrist straps when touching internal components.
	\item Avoid abrupt system shutdowns during model training or database write operations, as this can corrupt files.
	\item Ensure cables are not under tension or twisted while the system is powered.
\end{itemize}

\subsection{Data and Software Safety}
\begin{itemize}
	\item Always back up the \textbf{SQLite database} and configuration files before major software updates.
	\item Use secure file permissions and avoid storing personal or sensitive vehicle data in plain text.
	\item Do not interrupt the model inference or OCR process during execution, as it may cause temporary resource lockups.
	\item Regularly check the system logs to monitor for overheating, low power warnings, or database errors.
\end{itemize}

\subsection{General Precautions}
\begin{itemize}
	\item Operate only under adult supervision in lab or test environments.
	\item Do not attempt to disassemble the Jetson module or heat sink assembly.
	\item Keep liquids, metallic objects, and magnetic fields away from the hardware setup.
	\item Ensure compliance with local data privacy and recording laws when using the camera in public areas.
\end{itemize}

% --------------------------------------------------
% Section 2: Overview
% --------------------------------------------------
\section{Overview}

\subsection{Hardware Description (Operator’s View)}
The setup includes a \textbf{Seeed Studio reComputer J1010} (Jetson Nano 4 GB) enclosed in an aluminum case.  
It is connected to a \textbf{Microsoft LifeCam Show USB} camera, an HDMI display, and powered by a 5 V / 4 A adapter.  
This configuration allows continuous video capture and AI inference directly at the edge.

\subsection{Preparation}
\begin{itemize}
	\item Place the unit on a flat, ventilated surface.
	\item Connect the USB camera, HDMI monitor, and power adapter.
	\item Insert the microSD card flashed with \textbf{JetPack 5.1.2 (Ubuntu 20.04)}.
	\item Power on the device and log in to the Ubuntu desktop.
	\item Launch the detection program via terminal or GUI.
\end{itemize}

% --------------------------------------------------
% TikZ Diagram – Hardware Connection
% --------------------------------------------------
\begin{figure}[h!]
	\centering
	\begin{tikzpicture}[node distance=2.8cm,>=stealth,thick]
		
		\tikzstyle{block}=[rectangle, draw=blue!50!black, fill=blue!10,
		rounded corners, align=center, text width=3.5cm, minimum height=1cm]
		
		\node[block] (camera) {\faCamera\\Camera\\(Microsoft LifeCam)};
		\node[block, right=of camera] (nano) {\faMicrochip\\Jetson Nano\\(reComputer J1010)};
		\node[block, right=of nano] (display) {\faDesktop\\HDMI Display};
		\node[block, below=of nano] (db) {\faDatabase\\SQLite Database};
		
		\draw[->, thick] (camera) -- node[above]{USB 2.0} (nano);
		\draw[->, thick] (nano) -- node[above]{HDMI Output} (display);
		\draw[->, thick] (nano) -- node[left]{Data Logging} (db);
	\end{tikzpicture}
	\caption{Hardware connection overview for License Plate Detection System}
	\label{fig:hardware_connection}
\end{figure}

% --------------------------------------------------
% Section 3: Operation
% --------------------------------------------------
\section{Operation – List of Functions}

\begin{itemize}
	\item \textbf{Function 1: Live Detection} — Captures real-time video and detects license plates with YOLOv8.
	\item \textbf{Function 2: OCR Recognition} — Extracts text from detected plates via EasyOCR.
	\item \textbf{Function 3: Database Logging} — Stores results in SQLite with timestamps.
	\item \textbf{Function 4: Exception Handling} — Skips billing for registered emergency vehicles.
	\item \textbf{Function 5: Reporting \& Display} — Shows detection results and logs on screen.
\end{itemize}

% --------------------------------------------------
% Section 4: Troubleshooting (Keep Box)
% --------------------------------------------------
\begin{tcolorbox}[colback=red!5!white,colframe=red!60!black,title={\faTools\hspace{0.5em}Quick Troubleshooting Hint }]
	\begin{itemize}
		\item \textbf{No Camera Feed:} Check USB connection; test with \texttt{ls /dev/video*}.
		\item \textbf{Low FPS:} Lower resolution or use the YOLOv8-Nano model.
		\item \textbf{Model Error:} Verify weights directory and filenames.
		\item \textbf{Database Locked:} Close previous app sessions and retry.
		\item \textbf{Overheating:} Ensure ventilation or add a small cooling fan.
	\end{itemize}
\end{tcolorbox}

% --------------------------------------------------
% Section 5: Specifications
% --------------------------------------------------
\section{Specifications}

\begin{table}[h!]
	\centering
	\renewcommand{\arraystretch}{1.4} % Increases row height for readability
	\setlength{\tabcolsep}{10pt} % Adjusts spacing between columns
	
	\begin{tcolorbox}[
		colback=gray!3!white,
		colframe=blue!60!black,
		title={\faMicrochip\hspace{0.5em}System Specifications},
		fonttitle=\bfseries,
		arc=3mm,
		boxrule=1pt,
		width=\textwidth-1cm, % prevents border overflow
		left=5mm, right=5mm, top=3mm, bottom=3mm
		]
		\newlength{\specwidth}
		\setlength{\specwidth}{\dimexpr\textwidth-0.5cm\relax}
\begin{tabularx}{\specwidth}{|>{\raggedright\arraybackslash}p{4cm}|X|}
			\hline
			\rowcolor{blue!10}
			\textbf{Category} & \textbf{Details} \\ \hline
			
			\textbf{Device} & Seeed Studio reComputer J1010 with NVIDIA Jetson Nano (4 GB LPDDR4) \\ \hline
			\textbf{Operating System} & JetPack 5.1.2 (Ubuntu 20.04 LTS) \\ \hline
			\textbf{Processor (CPU)} & Quad-core ARM Cortex-A57 @ 1.43 GHz \\ \hline
			\textbf{Graphics (GPU)} & 128-core NVIDIA Maxwell architecture, CUDA-capable \\ \hline
			\textbf{Memory} & 4 GB LPDDR4 (25.6 GB/s bandwidth) \\ \hline
			\textbf{Storage} & 64 GB Class 10 microSD card (expandable) \\ \hline
			\textbf{Camera Module} & Microsoft LifeCam Show (USB 2.0, 720p) \\ \hline
			\textbf{Connectivity} & USB 2.0, HDMI output, Gigabit Ethernet \\ \hline
			\textbf{Power Supply} & 5 V / 4 A via USB-C connector \\ \hline
			\textbf{Software Stack} & YOLOv8 (Ultralytics), EasyOCR, SQLite3, OpenCV, NumPy, Python 3.8 \\ \hline
			\textbf{Average Performance} & 13–15 FPS detection rate with YOLOv8-Nano (real-time) \\ \hline
			\textbf{Use Case} & Real-time vehicle monitoring, parking automation, and access control \\ \hline
			
		\end{tabularx}
	\end{tcolorbox}
	
	\caption{Detailed hardware and software specifications of the Jetson Nano License Plate Detection System.}
	\label{tab:specifications}
\end{table}